\section{مدول‌های نوتری و آرتینی}

% Source: sources/lectures/10 Notherian and Artinian modules.pdf | PDF page: 1 | item: outside-course remark on real numbers
\begin{frame}[t]{اعداد حقیقی به‌عنوان مدول}
    \label{lec10:real-modules}
    \begin{remark}[خارج از درس]
        $\mathbb R$ به‌عنوان $\mathbb Z$-مدول آزاد نیست؛
        زیرا $\mathbb Q$ به‌عنوان $\mathbb Z$-مدول آزاد نیست و
        زیرگروه هر گروه آبلی آزاد، آزاد است.
        در جزوه به قضیهٔ ددکیند یا
        \lr{Nielsen--Schreier} اشاره شده است.
    \end{remark}
    \begin{remark}
        $\mathbb R$ به‌عنوان $\mathbb Q$-مدول آزاد است
        و پایه‌ای موسوم به پایهٔ \lr{Hamel} دارد.
    \end{remark}
\end{frame}

% Source: sources/lectures/10 Notherian and Artinian modules.pdf | PDF page: 1 | item: finite free module is isomorphic to R^n; source proof ends in ellipsis
\begin{frame}[t]{مدول آزاد با پایهٔ متناهی}
    \label{lec10:finite-free}
    \begin{lemma}
        اگر $M$ یک $R$-مدول آزاد با پایهٔ
        $\{m_1,\ldots,m_n\}$ باشد، آنگاه
        \[
            M \modulecong R^n.
        \]
    \end{lemma}
    \textbf{طرح برهان جزوه.} نگاشت زیر معرفی می‌شود:
    \[
        \begin{aligned}
            \psi:R^n&\longrightarrow M,\\
            (\alpha_1,\ldots,\alpha_n)&\longmapsto
                \sum_{i=1}^n\alpha_i m_i.
        \end{aligned}
    \]
    ادامهٔ بررسی این نگاشت در منبع با سه‌نقطه واگذار شده است.
\end{frame}

% Source: sources/lectures/10 Notherian and Artinian modules.pdf | PDF pages: 1-2 | item: every basis of a finitely generated free module is finite; finite supports
\begin{frame}[t]{متناهی بودن هر پایه: مولدها}
    \label{lec10:finite-basis-supports}
    \begin{theorem}
        اگر $M$ یک $R$-مدول آزاد با تولید متناهی باشد،
        هر پایهٔ $M$ متناهی است.
    \end{theorem}
    \textbf{اثبات.} یک پایهٔ دلخواه $E=\{e_i\}_{i\in I}$ و
    مولدهای متناهی $M$ را در نظر بگیرید:
    \[
        M=\langle m_1,\ldots,m_n\rangle.
    \]
    هر مولد، ترکیبی خطی و متناهی از اعضای $E$ است:
    \[
        m_j=\sum_{\ell=1}^{k_j}r_{j\ell}e_{i_{j\ell}},
        \qquad 1\leq j\leq n.
    \]
    مجموعهٔ تمام اعضای پایه که در این نمایش‌ها ظاهر می‌شوند
    را $T$ بنامید. مجموعهٔ $T$ متناهی است.
\end{frame}

% Source: sources/lectures/10 Notherian and Artinian modules.pdf | PDF page: 2 | item: T generates M and linear independence forces E=T
\begin{frame}[t]{متناهی بودن هر پایه: نتیجه}
    \label{lec10:finite-basis-conclusion}
    چون همهٔ مولدهای $m_1,\ldots,m_n$ در $\langle T\rangle$ قرار دارند،
    \[
        M=\langle T\rangle,\qquad T\subseteq E.
    \]
    از استقلال خطی $E$، مجموعهٔ $T$ نیز مستقل خطی است؛
    پس $T$ یک پایهٔ متناهی برای $M$ است.

    اگر عضوی از $E$ بیرون $T$ باشد، چون $T$ مولد $M$ است،
    آن عضو ترکیبی خطی از اعضای $T$ خواهد بود.
    این با استقلال خطی $E$ تناقض دارد. بنابراین
    \[
        E=T,
    \]
    و پایهٔ دلخواه $E$ متناهی است.
    \hfill\(\qed\)
\end{frame}

% Source: sources/lectures/10 Notherian and Artinian modules.pdf | PDF pages: 2-3 | item: equal cardinalities of two bases; commutative unital hypothesis
% AMBIGUITY[L10-Q01]: nonzero-ring hypothesis is pending; do not silently add it.
\begin{frame}[t]{برابری تعداد اعضای دو پایه}
    \label{lec10:equal-bases}
    \begin{tblock}{پیش‌نویس؛ پرسش \lr{L10-Q01}}
        فرض ناصفر بودن حلقه در صورت زیر هنوز نیازمند تأیید است.
    \end{tblock}
    \textbf{صورت جزوه.} فرض کنید $R$ حلقه‌ای جابجایی و یکدار و
    $M$ یک $R$-مدول آزاد با تولید متناهی باشد.
    در این صورت هر دو پایهٔ $M$ تعداد برابری عضو دارند.

    \textbf{شروع برهان.} برای دو پایهٔ
    $S_1=\{e_1,\ldots,e_n\}$ و
    $S_2=\{e'_1,\ldots,e'_m\}$ داریم:
    \[
        M\cong R^n,\qquad M\cong R^m
        \quad\Longrightarrow\quad R^m\modulecong R^n.
    \]
    برای رسیدن به تناقض، مانند جزوه فرض کنید $m<n$.
\end{frame}

% Source: sources/lectures/10 Notherian and Artinian modules.pdf | PDF page: 3 | item: inverse matrices, square padding, determinant contradiction
% AMBIGUITY[L10-Q01]: 0=1 is a contradiction only once 1_R != 0_R is confirmed.
\begin{frame}[t]{برهان دترمینانی برابری دو پایه}
    \label{lec10:equal-bases-determinant}
    \textbf{ادامهٔ پیش‌نویس \lr{L10-Q01}.}
    یکریختی‌های معکوس $\varphi:R^m\to R^n$ و $\psi:R^n\to R^m$
    با ماتریس‌های $A_{n\times m}$ و $B_{m\times n}$ در جزوه به روابط زیر می‌رسند:
    \[
        AB=I_n,\qquad BA=I_m.
    \]
    با افزودن بلوک‌های صفر، دو ماتریس مربعیِ $n\times n$ بسازید:
    \[
        \bar A=\begin{bmatrix}A&0\end{bmatrix},
        \qquad
        \bar B=\begin{bmatrix}B\\0\end{bmatrix}.
    \]
    \[
        \bar A\bar B=I_n,\qquad
        \bar B\bar A=
        \begin{bmatrix}I_m&0\\0&0\end{bmatrix}.
    \]
    پس از برابری دترمینان دو حاصل‌ضرب، $1_R=0_R$ حاصل می‌شود.
    نتیجه‌گیریِ تناقض، منتظر تعیین فرض ناصفر بودن $R$ است.
\end{frame}

% Source: sources/lectures/10 Notherian and Artinian modules.pdf | PDF page: 4 | item: definitions of Noetherian and Artinian modules
\begin{frame}[t]{شرط‌های زنجیری}
    \label{lec10:chain-conditions}
    \begin{definition}
        $R$-مدول $M$ را \textbf{نوتری} گوییم هرگاه هر زنجیر صعودی
        از $R$-زیرمدول‌های آن متوقف شود.
    \end{definition}
    \[
        M_1\subseteq M_2\subseteq M_3\subseteq\cdots
    \]
    \begin{definition}
        $R$-مدول $M$ را \textbf{آرتینی} گوییم هرگاه هر زنجیر نزولی
        از $R$-زیرمدول‌های آن متوقف شود.
    \end{definition}
    \[
        M_1\supseteq M_2\supseteq M_3\supseteq\cdots
    \]
\end{frame}

% Source: sources/lectures/10 Notherian and Artinian modules.pdf | PDF page: 4 | item: Z is Noetherian but not Artinian; source why-prompt retained
\begin{frame}[t]{اعداد صحیح: نوتری ولی ناآرتینی}
    \label{lec10:integers}
    \begin{example}
        $\mathbb Z$ به‌عنوان $\mathbb Z$-مدول نوتری است.
        همهٔ زیرمدول‌های آن به شکل $\langle m\rangle$ هستند و
        \[
            n\mid m
            \quad\Longleftrightarrow\quad
            \langle m\rangle\subseteq\langle n\rangle.
        \]
        \textbf{تمرین جزوه:} چرا هر زنجیر صعودی از این زیرمدول‌ها
        متوقف می‌شود؟
    \end{example}
    \begin{example}
        $\mathbb Z$ به‌عنوان $\mathbb Z$-مدول آرتینی نیست؛ زیرا
        \[
            \langle2\rangle\supsetneq\langle4\rangle
            \supsetneq\langle8\rangle\supsetneq\cdots
        \]
        هیچ‌گاه متوقف نمی‌شود.
    \end{example}
\end{frame}

% Source: sources/lectures/10 Notherian and Artinian modules.pdf | PDF page: 4 | item: finite abelian groups, Q, finite-dimensional vector spaces
\begin{frame}[t]{چند مثال دیگر}
    \label{lec10:chain-examples}
    \begin{example}
        هر گروه آبلی متناهی، به‌عنوان $\mathbb Z$-مدول،
        هم نوتری و هم آرتینی است.
    \end{example}
    \begin{example}
        $\mathbb Q$ به‌عنوان $\mathbb Z$-مدول نه نوتری است و نه آرتینی:
        \[
            \langle2\rangle\supsetneq\langle4\rangle
            \supsetneq\langle8\rangle\supsetneq\cdots,
        \]
        \[
            \left\langle\tfrac12\right\rangle
            \subsetneq\left\langle\tfrac14\right\rangle
            \subsetneq\left\langle\tfrac18\right\rangle
            \subsetneq\cdots.
        \]
    \end{example}
    \begin{example}
        هر فضای برداری با بعد متناهی، بر میدانِ ضرایب خود،
        هم نوتری و هم آرتینی است.
    \end{example}
\end{frame}

% Source: sources/lectures/10 Notherian and Artinian modules.pdf | PDF page: 5 | item: definition of the Prüfer group and quotient notation
\begin{frame}[t]{گروه پروفر}
    \label{lec10:prufer-definition}
    \begin{definition}
        فرض کنید $p$ عددی اول باشد و قرار دهید:
        \[
            M=\left\{\frac{a}{p^n}\ \middle|\ %
                a\in\mathbb Z,\ n\in\mathbb N\right\}.
        \]
        $(M,+)$ یک گروه آبلی است و
        \[
            (\mathbb Z,+)\subsetneq M\subsetneq(\mathbb Q,+).
        \]
        گروه خارج‌قسمتی $M/\mathbb Z$ را \textbf{گروه پروفر} می‌نامیم و
        با $\mathbb Z_{p^\infty}$ نمایش می‌دهیم:
        \[
            \frac{a}{p^n}+\mathbb Z
            =\overline{\frac{a}{p^n}}.
        \]
    \end{definition}
\end{frame}

% Source: sources/lectures/10 Notherian and Artinian modules.pdf | PDF page: 5 | item: G_n and claims 1-3; source why-prompts retained
\begin{frame}[t]{یک زنجیر از زیرگروه‌های پروفر}
    \label{lec10:prufer-chain}
    برای $n\geq1$ تعریف کنید:
    \[
        G_n=\left\{\overline{\frac{a}{p^n}}
            \in\mathbb Z_{p^\infty}\ \middle|\ a\in\mathbb Z\right\}.
    \]
    \begin{enumerate}
        \item هر $G_n$ یک $\mathbb Z$-زیرمدولِ $\mathbb Z_{p^\infty}$ است.
        \item برای هر $n$ داریم $G_n\subsetneq G_{n+1}$.
        \item $\displaystyle\mathbb Z_{p^\infty}=\bigcup_{n=1}^{\infty}G_n$.
    \end{enumerate}
    \textbf{تمرین‌های جزوه:} زیرمدول بودن، شمول و اکید بودن شمول
    را بررسی کنید.
\end{frame}

% Source: sources/lectures/10 Notherian and Artinian modules.pdf | PDF pages: 5-6 | item: Bezout argument from a reduced representative to G_n contained in N
\begin{frame}[t]{مولد یک زیرگروه پروفر}
    \label{lec10:prufer-bezout}
    فرض کنید $N\leq\mathbb Z_{p^\infty}$ و
    \[
        \overline{\frac{a}{p^n}}\in N,\qquad p\nmid a.
    \]
    چون $(a,p)=1$، داریم $(a,p^n)=1$؛ پس
    \[
        \exists r,s\in\mathbb Z:\qquad ra+sp^n=1.
    \]
    بنابراین
    \[
        r\,\overline{\frac{a}{p^n}}
        =\overline{\frac{ra}{p^n}}
        =\overline{\frac{1-sp^n}{p^n}}
        =\overline{\frac1{p^n}}-\overline{s}
        =\overline{\frac1{p^n}}\in N,
    \]
    زیرا $\overline{s}=0$ در $M/\mathbb Z$. در نتیجه
    \[
        G_n\subseteq N.
    \]
\end{frame}

% Source: sources/lectures/10 Notherian and Artinian modules.pdf | PDF pages: 5-6 | item: claim 4 conflicts with the whole-group case on the following page
% AMBIGUITY[L10-Q02]: proper subgroup and reduced-denominator scope require user confirmation.
\begin{frame}[t]{طبقه‌بندی زیرگروه‌های پروفر: پرسش باز}
    \label{lec10:prufer-classification-question}
    \begin{tblock}{پیش‌نویس؛ پرسش \lr{L10-Q02}}
        ادعای چهارم صفحهٔ ۵ و حالت دوم صفحهٔ ۶ به تعیین دامنهٔ ادعا نیاز دارند.
    \end{tblock}
    \textbf{صورت ادعای چهارم در منبع:}
    \[
        0\ne N\leq\mathbb Z_{p^\infty}
        \quad\Longrightarrow\quad
        \exists n:\ N=G_n.
    \]
    \textbf{ولی در صفحهٔ بعد آمده است:}
    اگر بزرگ‌ترین توان مخرج وجود نداشته باشد،
    \[
        N=\mathbb Z_{p^\infty}.
    \]
    قید «زیرگروه سره» در ادعا و استفاده از نمایش‌های تحویل‌یافته
    در انتخاب توان مخرج، برای تأیید در فایل پرسش‌ها ثبت شده‌اند.
\end{frame}

% Source: sources/lectures/10 Notherian and Artinian modules.pdf | PDF page: 6 | item: maximal exponent proof, unbounded case, and chain conclusion
% AMBIGUITY[L10-Q02]: source proof outline quoted provisionally; classification and Artinian proof not finalized.
\begin{frame}[t]{طرح برهان طبقه‌بندی در جزوه}
    \label{lec10:prufer-classification-outline}
    \textbf{ادامهٔ پیش‌نویس \lr{L10-Q02}.}
    \begin{itemize}
        \item اگر بزرگ‌ترین توان مخرج، $n$ باشد، برهان قبل $G_n\subseteq N$
              را می‌دهد. جزوه برای $N\ne G_n$ عضوی با توان مخرج $m>n$
              و $(a,p)=1$ در نظر می‌گیرد و به تناقض می‌رسد.
        \item اگر چنین بیشینه‌ای وجود نداشته باشد،
              جزوه نتیجه می‌گیرد $N=\mathbb Z_{p^\infty}$.
    \end{itemize}
    نتیجهٔ نوشته‌شده در منبع، با نماد $G_0$ برای زیرگروه صفر:
    \[
        G_0\subsetneq G_1\subsetneq G_2\subsetneq G_3\subsetneq\cdots
    \]
    و آرتینی بودن و نوتری نبودن $\mathbb Z_{p^\infty}$ است.

    \textbf{وضعیت:} نوتری نبودن از زنجیر بالا روشن است؛
    صورت دقیق طبقه‌بندی و تکمیل برهان آرتینی بودن، منتظر پاسخ پرسش است.
\end{frame}

% Source: sources/lectures/10 Notherian and Artinian modules.pdf | PDF page: 6 | item: three equivalent Noetherian conditions for a unital module
\begin{frame}[t]{معیارهای نوتری بودن}
    \label{lec10:noetherian-criteria}
    \begin{theorem}
        اگر $M$ یک $R$-مدول یکانی باشد، گزاره‌های زیر معادل‌اند:
        \begin{enumerate}
            \item $M$ یک $R$-مدول نوتری است.
            \item هر زیرمدول $M$ با تولید متناهی است.
            \item هر مجموعهٔ ناتهی از زیرمدول‌های $M$ یک عضو ماکسیمال دارد.
        \end{enumerate}
    \end{theorem}
    در ادامه، گذارهای زیر بررسی می‌شوند:
    \[
        (1)\Longrightarrow(2)\Longrightarrow(3)\Longrightarrow(1).
    \]
\end{frame}

% Source: sources/lectures/10 Notherian and Artinian modules.pdf | PDF page: 7 | item: (1) implies (2), iterative construction of generators
\begin{frame}[t]{از شرط زنجیری به تولید متناهی}
    \label{lec10:noetherian-finite-generation}
    \textbf{اثبات \((1)\Longrightarrow(2)\).}
    فرض کنید $N\leq M$. برای $N=0$ حکم روشن است.
    در غیر این صورت، $0\ne x_1\in N$ انتخاب کنید.

    اگر $N=\langle x_1\rangle$، کار تمام است؛ وگرنه
    \[
        x_2\in N\setminus\langle x_1\rangle.
    \]
    اگر $N=\langle x_1,x_2\rangle$، کار تمام است؛ وگرنه
    \[
        x_3\in N\setminus\langle x_1,x_2\rangle.
    \]
    ادامهٔ فرایند، زنجیر اکید زیر را می‌سازد:
    \[
        \langle0\rangle\subsetneq\langle x_1\rangle
        \subsetneq\langle x_1,x_2\rangle\subsetneq\cdots.
    \]
    نوتری بودن $M$ فرایند را متوقف می‌کند؛ پس $N$ با تولید متناهی است.
\end{frame}

% Source: sources/lectures/10 Notherian and Artinian modules.pdf | PDF page: 7 | item: (2) implies (3), strict chain, union, finitely many generators
\begin{frame}[t]{از تولید متناهی به شرط ماکسیمال}
    \label{lec10:noetherian-maximal-family}
    \textbf{شروع برهان \((2)\Longrightarrow(3)\).}
    یک مجموعهٔ ناتهی از زیرمدول‌های $M$ را در نظر بگیرید.
    اگر عضو ماکسیمال نداشته باشد، با انتخاب پیاپی اعضای آن
    زنجیر اکید زیر ساخته می‌شود:
    \[
        M_1\subsetneq M_2\subsetneq M_3\subsetneq\cdots.
    \]
    قرار دهید
    \[
        N=\bigcup_{i=1}^{\infty}M_i.
    \]
    $N$ زیرمدول $M$ و کران بالای همهٔ $M_i$ها است.
    طبق فرض، مولدهای متناهی $n_1,\ldots,n_k$ وجود دارند که
    \[
        N=\langle n_1,\ldots,n_k\rangle.
    \]
\end{frame}

% Source: sources/lectures/10 Notherian and Artinian modules.pdf | PDF page: 7 | item: final chain index in the (2) implies (3) proof
% AMBIGUITY[L10-Q03]: proposed chain index K=max{j_1,...,j_k} is not yet approved.
\begin{frame}[t]{اندیس توقف زنجیر: پرسش باز}
    \label{lec10:noetherian-maximal-index}
    \begin{tblock}{پیش‌نویس؛ پرسش \lr{L10-Q03}}
        اندیس سطر آخر برهان در جزوه نیازمند تعیین دقیق است.
    \end{tblock}
    هر مولد $n_i$ در یکی از اعضای زنجیر قرار دارد.
    در منبع، پس از این گام $N=M_k$ نوشته شده است.

    \textbf{تکمیل پیشنهادی، در انتظار تأیید:}
    \[
        n_i\in M_{j_i}\quad(1\leq i\leq k),\qquad
        K=\max\{j_1,\ldots,j_k\}.
    \]
    \[
        N=\langle n_1,\ldots,n_k\rangle\subseteq M_K\subseteq N.
    \]
    در این صورت $N=M_K$ و
    $M_K\subsetneq M_{K+1}\subseteq N$ تناقض می‌دهد.
    این انتخاب اندیس و تکمیل برهان هنوز قطعی نشده است.
\end{frame}

% Source: sources/lectures/10 Notherian and Artinian modules.pdf | PDF page: 8 | item: (3) implies (1), apply condition (3) to the chain itself
\begin{frame}[t]{از شرط ماکسیمال به توقف زنجیر}
    \label{lec10:noetherian-acc}
    \textbf{اثبات \((3)\Longrightarrow(1)\).}
    فرض کنید
    \[
        M_1\subseteq M_2\subseteq M_3\subseteq\cdots
    \]
    یک زنجیر صعودی از زیرمدول‌های $M$ باشد.

    مانند جزوه، شرط (۳) را به خودِ مجموعهٔ ناتهی
    $\{M_i\}_{i=1}^{\infty}$ اعمال کنید.
    این مجموعه یک عضو ماکسیمال، مثلاً $M_k$، دارد.

    برای هر $i\geq k$، از صعودی بودن زنجیر $M_k\subseteq M_i$؛
    از ماکسیمال بودن $M_k$ نیز $M_i=M_k$. بنابراین
    \[
        M_k=M_{k+1}=M_{k+2}=\cdots,
    \]
    و زنجیر متوقف می‌شود.
\end{frame}

% Source: sources/lectures/10 Notherian and Artinian modules.pdf | PDF page: 8 | item: finite generation by finite subfamilies; equivalence is a source exercise
\begin{frame}[t]{تولید متناهی با خانوادهٔ زیرمدول‌ها}
    \label{lec10:finite-generation-family}
    \begin{definition}
        $R$-مدول $M$ را \textbf{متناهیاً تولیدشده} گوییم هرگاه
        برای هر خانوادهٔ $\{M_i\}_{i\in I}$ از $R$-زیرمدول‌های آن،
        از
        \[
            M=\sum_{i\in I}M_i
        \]
        بتوان نتیجه گرفت که زیرمجموعه‌ای متناهی $J\subseteq I$ وجود دارد با
        \[
            M=\sum_{i\in J}M_i.
        \]
    \end{definition}
    \textbf{تمرین جزوه:} معادل بودن این تعریف با تعریف قبلیِ
    تولید متناهی را بررسی کنید.
\end{frame}

% Source: sources/lectures/10 Notherian and Artinian modules.pdf | PDF page: 8 | item: definition of finite cogeneration
\begin{frame}[t]{هم‌تولید متناهی}
    \label{lec10:finite-cogeneration}
    \begin{definition}
        $R$-مدول $M$ را \textbf{متناهیاً هم‌تولیدشده} گوییم هرگاه:

        برای هر خانوادهٔ $\{M_i\}_{i\in I}$ از $R$-زیرمدول‌های آن،
        از
        \[
            \bigcap_{i\in I}M_i=0
        \]
        بتوان نتیجه گرفت که زیرمجموعه‌ای متناهی $J\subseteq I$ وجود دارد با
        \[
            \bigcap_{i\in J}M_i=0.
        \]
    \end{definition}
\end{frame}
